\chapter*{On the Recurrence of \texorpdfstring{$2/3$}{2/3}: A Reply to Austin's Algebraic Note}
\addcontentsline{toc}{chapter}{Reply to Austin's Algebraic Note on $2/3$}

\begin{flushright}\small
Johann Pascher \\
in reply to Peter Austin, \emph{A Purely Algebraic Short Note on the
Repeated Appearance of $2/3$} (May 2026) \\
for the IPI correspondence group --- May 2026
\end{flushright}

\section{Purpose and stance}

Peter, thank you for the note. It is careful, the algebra is correct, and
it reconstructs the FFGFT lepton sector accurately from the public
summaries. What follows is not a defence and not a request for validation:
it is the set of internal derivations your note effectively asks for. Where
your reconstruction stops at the public numbers, FFGFT already carries the
structure one level deeper, and in two places your note independently
arrives at conclusions that the framework reached on its own. I will use
your section numbering throughout.

For the record, the canonical charged-lepton coefficients are
\[
  (r_e, r_\mu, r_\tau) = \left(\tfrac{4}{3},\, \tfrac{16}{5},\, \tfrac{25}{9}\right),
  \qquad
  (p_e, p_\mu, p_\tau) = \left(\tfrac{3}{2},\, 1,\, \tfrac{2}{3}\right),
\]
with $m_i = r_i\,\xipar^{p_i}\,v$ and $\xipar = 4/30000$. Your note uses
exactly these; they are current.

\section{The $\epsilon$-deformation discriminator (your \S6)}

I accept the discriminator without reservation. Coincidence at $D=3$ is not
identity of mechanism, and the continuation $D = 3+\epsilon$ is the right
instrument to separate the roads. This is worth stating plainly because
FFGFT reaches the \emph{same} separation from the inside, by a different
route (see \S5 below): the exponent-ladder value of the Koide quantity and
the equal-projector condition $A=B$ do not coincide exactly --- they meet
\emph{near} $2/3$ and part company, precisely as your $\epsilon$-analysis
predicts. We are not in disagreement on the method; we agree on the
conclusion.

\section{The ladder is geometric, not arithmetic (your \S11)}

Your \S11 is correct that $\left(\tfrac{3}{2}, 1, \tfrac{2}{3}\right)$ is
not a uniform $1/3$ ladder: the gaps are
\[
  p_e - p_\mu = \tfrac{1}{2}, \qquad p_\mu - p_\tau = \tfrac{1}{3}.
\]
The resolution is that the ladder is not arithmetic at all. It is a
\emph{geometric} sequence with ratio $q = 2/3$:
\[
  p_{n+1} = \tfrac{2}{3}\,p_n,
  \qquad
  \tfrac{3}{2} \xrightarrow{\times 2/3} 1 \xrightarrow{\times 2/3} \tfrac{2}{3}.
\]
The unequal differences $1/2$ and $1/3$ are then \emph{consequences} of a
constant ratio, not independent rules to be posted. So the electron does
not need its own ad-hoc rule; it is the base rung of the sequence. Two
equivalent readings of the same fact:

\begin{itemize}
  \item \textbf{Ratio reading.} $q = 2/3$ is the inverse of the
        tetrahedral geometry factor $3/2$ (four vertices over three edges
        per vertex). The same number that orders the ladder is, numerically,
        the Koide value --- which is exactly why the two appear linked.
  \item \textbf{Complement reading (your road).} At the $\mu\!\to\!\tau$
        rung the step is $1 - 1/D = 2/3$ with $D = 3$, so
        $p_\tau = 1 - 1/D$. This places FFGFT on your ``complement of
        step'' road, as you correctly identified, and \emph{not} on the
        equal-projector road $2/D$.
\end{itemize}

In the harmonic picture the rungs sit on a logarithmic torus in steps of a
musical third, $\Delta p = 1/3$. The electron at $p_e = 3/2 = 1 + 1/(D-1)$
is a \emph{half-integer} rung --- a geometric midpoint between winding
levels (the spin-winding contribution), not a full winding step. That is
the structural origin of the $1/2$ gap.

\section{Topological $D=3$ versus the fractal dimension (your \S11--\S12)}

Your \S12 asks whether the $1/3$ step is tied to the integer $D=3$ or to
$D_f = 3-\xi$, the latter giving a downward shift $2/3 - \xi/9$. The
framework's answer is specific, and it dissolves the worry for the Koide
case in particular.

First, there are \emph{two} derived fractal dimensions, not one:
\[
  D_f^{\,\text{space}} = 3 - \xi \approx 2.9998667
  \quad\text{(local metric; GRB time-of-flight)},
\]
\[
  D_f^{\,\text{eff}} \approx 2.973
  \quad\text{(cumulative scale flow, Planck to electroweak)}.
\]
Both are \emph{derived}, not assumed. The exponent step is tied to the
\emph{integer, topological} torus dimension $D=3$; the metric $3-\xi$ enters
elsewhere (dispersion), and the effective $2.973$ enters only the absolute
normalisation through $K_{\text{frak}} = (D_f^{\,\text{eff}}/3)^{D_f^{\,\text{eff}}/2}
= 1 - 100\,\xi$.

\begin{keyresult}[Why no $\xi/9$ correction reaches Koide]
The Koide quantity $Q$ is a pure mass \emph{ratio}. Every scale-dependent
factor --- the vacuum value $v$, the fractal correction $K_{\text{frak}}$,
and any $D_f$-level shift --- is common to all three leptons and cancels in
$Q$. Fractal corrections attach to \emph{absolute} scales, never to ratios.
Your \S12 shift $2/3 - \xi/9$ would be a real effect on an absolute
prediction; it simply does not enter $Q$, because $Q$ never sees the
absolute scale.
\end{keyresult}

So for Koide the choice is settled in favour of your option (a): the step
is on topological $D=3$, and $D_f$ does not perturb the ratio.

\section{Forced or merely fitted? An honest answer (your \S10, \S13)}

This is the central question of your note, and FFGFT's internal answer is
already on record --- and it is more interesting than a clean ``forced.''

\subsection*{What is structurally forced}

The mass \emph{ratios} are not fitted. Two points:
\begin{itemize}
  \item $m_\mu/m_e$ follows two independent ways that agree:
        geometrically from the ladder, $\frac{12}{5}\,\xi^{-1/2} \approx 207.8$,
        and fractally from $(D_f^{\,\text{eff}}/3)^{D_f^{\,\text{eff}}/2}\!\cdot f(\xi)
        \approx 206.8$. The agreement fixes $K_{\text{frak}}$ uniquely --- a
        consistency test, not a fit.
  \item In the compact representation the coefficients are not free
        rationals: $\tfrac{2}{3} = 3\sqrt{3}/(2\pi\alpha^{1/2})$ and
        $\tfrac{8}{5} = 9/(4\pi\alpha)$, built so that $\alpha$ cancels
        exactly. The rationals are $\alpha$-geometric invariants.
\end{itemize}

\subsection*{What the ladder does \emph{not} force}

It does not force the equal-projector condition $A=B$ exactly. Inserting the
FFGFT masses into the Koide formula gives
\[
  Q_{\text{FFGFT}} \approx 0.6677,
  \qquad
  \Delta_K = 3Q_{\text{FFGFT}} - 2 \approx +0.003,
\]
i.e.\ your residual diagnostic $\Delta_K = B/A - 1$ is small but nonzero:
$A = B$ holds to about $0.16\%$ in $Q$, not exactly. This is deliberate and
honest. The reason is structural, and it is your own point in disguise: the
Koide denominator $\big(\sum \sqrt{m_i}\big)^2$ generates cross-terms with
exponents
\[
  \xi^{5/4}, \quad \xi^{13/12}, \quad \xi^{5/6},
\]
none of which is a multiple of $1/3$. They lie \emph{outside} the FFGFT
structure space $\{p\} \subset \tfrac{1}{3}\mathbb{Z}$; they are artefacts of
the square-root composition, not torus modes. Exact $Q = 2/3$ would require
\[
  m_e + m_\mu + m_\tau = 4\big(\sqrt{m_e m_\mu} + \sqrt{m_e m_\tau} + \sqrt{m_\mu m_\tau}\big),
\]
which holds not at $\xipar = 4/30000$ but at $\xi_{\text{Koide}} \approx
1.372\times 10^{-4}$, about $2.9\%$ larger. In other words, the
equal-projector condition is an algebraic side-condition that is
\emph{independent} of the torus geometry --- which is exactly what your
$\epsilon$-deformation says it should be: the ladder road and the
equal-projector road are different mechanisms that meet near $2/3$ and
separate.

\begin{keypoint}[The honest position]
The ladder \emph{predicts} $Q$ to $0.16\%$ from $(r_i, p_i, \xi)$ with no
Koide input. The exact experimental $Q = 2/3$ then \emph{confirms} the
exponent structure from an independent direction; it is neither fitted into
FFGFT nor algebraically forced by it. The small residual is real and
located precisely where your cross-term analysis predicts.
\end{keypoint}

\section{Three roads and the map (your \S13)}

I agree with your framing: the Koide projector form, the FFGFT $\xi$-ladder,
and Paul's $24$-cell rank split are three different primitive objects, and a
relation between them must be exhibited by a map, not by a shared numeral.
The FFGFT-to-Koide map is explicit,
\[
  (r_i, p_i, \xi)\ \longmapsto\ x_i = \sqrt{r_i}\,\xi^{p_i/2}\,\sqrt{v}
  \ \longmapsto\ (A, B)\ \longmapsto\ Q,
\]
and your trigonometry-free target is, in our notation, identical to the
generated quantity
\[
  Q_{\text{FFGFT}} =
  \frac{\sum_i r_i\,\xi^{p_i}}{\big(\sum_i \sqrt{r_i}\,\xi^{p_i/2}\big)^2},
\]
since $\sqrt{v}$ cancels. We are computing the same object. The point of
disagreement, if any, is not the map but the claim of exactness: the map
sends the ladder to $Q \approx 0.6677$, not to $2/3$ on the nose, for the
reason in \S5.

\section{Closing}

So, briefly: your method is right; your placement of FFGFT on the
complement road is right; the ladder asymmetry resolves as a geometric
sequence of ratio $2/3$; the $D_f$ shift does not reach Koide because $Q$
is a ratio; and the ``forced versus fitted'' question has the honest answer
that the ladder predicts $Q$ to $0.16\%$ and the exact value is a
confirmation along an independent road --- which your own $\epsilon$-analysis
anticipates. If it is useful, you are welcome to post this alongside your
note; I think the two together make the cleaner statement.

\begin{flushright}\small
--- J.P.
\end{flushright}
